\subsection{半局部环}

命题 16. 设 A 是一个环。下列性质等价：

\begin{enumerate}
\def\labelenumi{(\alph{enumi})}
\item
  \(\mathbf{A}\) 的极大理想集合是有限的；
\item
  A 对其 Jacobson 根的商是有限个域的直合成。
\end{enumerate}

设 A 对其 Jacobson 根 \(\Re\) 的商是有限个域的直合成。那么 \(\Re\) 只有有限多个理想，\(a\ fortiori\) 只有有限多个极大理想。由于每个极大理想都包含 \(\Re\)（《代数》第 VIII 章 §6 第 2 节定义 2），A 的极大理想是 \(\Re\) 的极大理想在典范同态 A \(\rightarrow\) \(\Re\) 下的原像；因而它们个数有限。

反之，设 A 只有有限多个互不相同的极大理想 \(m_{1}, \ldots, m_{n}\)。诸 \(A/m_{i}\) 是域，且由 §1 第 2 节命题 5，典范映射 \(A \to \prod_{i=1}^{n} A/m_{i}\) 是满射；由于其核 \(\bigcap_{i=1}^{n} m_{i}\) 是 Jacobson 根 R（《代数》第 VIII 章 §6 第 2 节定义 2），A/R 同构于 \(\prod_{i=1}^{n} A/m_{i}\)。

定义 4. 若一个环满足命题 16 的等价条件 (a)、(b)，则称它为半局部环。

例. 每个局部环都是半局部的。半局部环的每个商环都是半局部的。半局部环的每个有限乘积都是半局部的。

* 若 A 是一个 Noether 半局部环，B 是一个作为 A-模为有限生成的 A-代数，则 B 是半局部的（第 IV 章 §2 第 5 节命题 9 的推论 3）。*

下述命题给出另一个例子，它推广了局部环 \(A_{p}\) 的构造：

命题 17. 设 A 是一个环，\(p_1, \ldots, p_n\) 是 A 的素理想。记 \(S = \bigcap_{i=1}^{n} (A - p_i) = A - \bigcup_{i=1}^{n} p_i\)。

\begin{enumerate}
\def\labelenumi{(\alph{enumi})}
\item
  环 \(\mathbf{S}^{-1}\mathbf{A}\) 是半局部的；若 \(\mathfrak{q}_1, \ldots, \mathfrak{q}_r\) 是诸 \(\mathfrak{p}_i\) 之中（按包含关系）互不相同的极大元，则 \(\mathbf{S}^{-1}\mathbf{A}\) 的极大理想是诸 \(\mathbf{S}^{-1}\mathfrak{q}_j\)（\(1 \leqslant j \leqslant r\)），且这些理想互不相同。
\item
  环 \(\mathbf{A}_{\mathfrak{p}_i}\) 典范同构于 \((\mathbf{S}^{-1}\mathbf{A})_{\mathbf{S}^{-1}\mathfrak{p}_i}\)（对 \(1 \leqslant i \leqslant n\)）。
\item
  若 \(\mathbf{A}\) 是一个整环，则 \(\mathbf{S}^{-1}\mathbf{A} = \bigcap_{i=1}^{n}\mathbf{A}_{p_i}\)（在 \(\mathbf{A}\) 的分式域中）。
\end{enumerate}

(a) A 的与 S 不相交的理想是含于诸 \(p_{i}\) 之并从而含于至少一个 \(p_{i}\) 之中的理想（§1 第 1 节命题 2）；因此诸 \(q_{j}\) 是与 S 不相交的理想集合的极大元；于是由 §2 第 5 节命题 11 (ii)，诸 \(S^{-1}q_{j}\) 是 \(S^{-1}A\) 的极大理想。

(b) 是 §2 第 5 节命题 11 (iii) 的特殊情形。

(c) 设 A 是一个整环。若 \(p_{i} \subset p_{k}\)，则 \(A_{p_{i}} \supset A_{p_{k}}\)；因此为证 (c)，可以设诸 \(p_{i}\) 两两不可比较。于是由 (a) 及第 3 节定理 1 的推论 4 得 \(S^{-1}A = \bigcap_{i=1}^{n} (S^{-1}A)_{S^{-1}p_{i}}\)；再凭 (b) 即得 (c)。

若 A 是整环，则 \(S^{-1}A\) 也是整环，于是命题 17 给出了一个不是局部环的直合成的半局部环的例子（参见~第 III 章 §2 第 13 节）。

推论. 设 A 是一个整环，\(p_{1}, \ldots, p_{n}\) 是 A 的素理想，其中任意两个按包含关系都不可比较。若在 A 的分式域中 \(A = \bigcap_{i=1}^{n} A_{p_{i}}\)，则 A 的极大理想是 \(p_{1}, \ldots, p_{n}\)。

取 \(S = \bigcap_{i=1}^{n} (A - p_i), S^{-1} A = A\)（由命题 17 (c)）；于是 \(S\) 的元素在 \(A\) 中可逆，且 \(S^{-1} p_i = p_i\) 对一切 \(i\) 成立。我们的断言随即凭命题 17 (a) 得出。

